\documentclass[11pt]{article}

\usepackage[utf8]{inputenc}
\usepackage[T1]{fontenc}
\usepackage[spanish]{babel}
\usepackage[margin=2.4cm]{geometry}
\usepackage{amsmath}
\usepackage{enumitem}
\usepackage{hyperref}
\hypersetup{colorlinks=true, urlcolor=blue, linkcolor=black}

\setlength{\parskip}{0.4em}

\title{\vspace{-1.5em}\textbf{Tarea 2 --- cmatch: threads y sincronización}\\[0.2em]
\large Sistemas Operativos}
\date{}

\begin{document}

% ----- Encabezado / portada compacta (se mantiene en una sola hoja para
%        respetar el límite de tres páginas que pide el enunciado) -----
\maketitle
\vspace{-3em}
\begin{center}
\begin{tabular}{ll}
Felipe Bustos López & \texttt{felipe.bustos@mail.udp.cl} \\
Cristóbal Sepúlveda Maldonado & \texttt{cristobal.sepulveda@mail.udp.cl} \\
\end{tabular}\\[0.4em]
Profesor: \textit{(indicar)} \hspace{2em} Mayo de 2026
\end{center}

\hrule
\vspace{0.6em}

% ---------------------------------------------------------------------------
\section{Introducción}

En esta tarea se construyó \texttt{cmatch}, un motor de emparejamiento
concurrente para torneos de \emph{tic-tac-toe} (gato). El sistema gestiona $N$
jugadores que buscan activamente formar partidas uno contra uno, compitiendo por
un número limitado de $K$ tableros. Cada jugador y cada tablero activo se modela
como un \emph{thread} independiente, por lo que el problema central no es la
lógica del juego --- que es trivial --- sino la coordinación segura de muchos
hilos que comparten recursos: la cola de espera, el puntaje ELO de los jugadores
y los tableros disponibles.

El objetivo de este informe es describir el diseño del programa, justificar los
mecanismos de sincronización empleados (únicamente la librería POSIX Threads,
\texttt{pthread}, según exige el enunciado) y comentar los mayores desafíos que
debimos resolver durante la implementación.

% ---------------------------------------------------------------------------
\section{Diseño del programa}

El programa se organiza en torno a dos tipos de hilos. Los \textbf{hilos de
jugador} ejecutan un ciclo de vida: ingresan a la cola de espera, buscan un
oponente válido, juegan una partida y, al terminar, deciden autónomamente
--- según la probabilidad \texttt{REENTER\_PROBABILITY} --- si reingresan a la
cola o finalizan su ejecución. Los \textbf{hilos de tablero} se crean por cada
partida y actúan como árbitros: coordinan los turnos, validan el estado del
juego y, al finalizar, actualizan el ELO y liberan el tablero.

El \textbf{emparejamiento} sigue las dos reglas pedidas. Un jugador en espera
busca entre los demás jugadores encolados aquel cuya diferencia de ELO no supere
$\Delta_{max}$ (\texttt{MAX\_ELO\_DIFF}); si hay varios candidatos válidos, se
elige al que lleva más tiempo esperando, lo que se resuelve asignando a cada
jugador un número de llegada creciente y prefiriendo el menor. El jugador que
encuentra rival y consigue un tablero libre actúa como ``anfitrión'': reclama a
ambos, crea la partida y lanza el hilo del tablero; el otro queda como
``invitado''. Así, ningún momento hay más de $K$ partidas simultáneas.

Al terminar una partida, el ELO de ambos jugadores se actualiza con el sistema
estándar de Elo,
\[
E_A = \frac{1}{1 + 10^{(\text{ELO}_B - \text{ELO}_A)/400}},
\qquad
\text{ELO}_{nuevo} = \text{ELO}_{ant} + K_{Elo}\,(S - E),
\]
donde $S \in \{1, 0.5, 0\}$ según el jugador gane, empate o pierda. La
configuración completa (cantidad de jugadores, tableros, factor $K_{Elo}$,
retardo entre jugadas, etc.) se lee de un archivo con formato \texttt{.env}.
Finalmente, el programa atiende la señal \texttt{SIGINT} para realizar un cierre
ordenado y guarda en disco un \emph{snapshot} con el estado de los jugadores,
que permite reanudar la simulación tal como quedó.

% ---------------------------------------------------------------------------
\section{Mecanismos de sincronización utilizados}

Decidimos usar \textbf{exclusivamente mutex y variables de condición} de
\texttt{pthread}, evitando incluso los semáforos POSIX (\texttt{<semaphore.h>}),
ya que el enunciado restringe la sincronización a la librería de hilos. Las
estructuras empleadas y su justificación son:

\begin{itemize}[leftmargin=1.4em, itemsep=0.2em]
  \item \textbf{Un mutex global (\texttt{g\_mutex})} que protege la ``tabla de
  jugadores'': la cola de espera, el estado de cada jugador, su ELO, sus
  estadísticas y el contador de tableros libres. Optamos por un \emph{bloqueo
  de grano grueso} porque estos datos se consultan y modifican de forma
  entrelazada durante el matchmaking; un único candado garantiza que, por
  ejemplo, jamás se lea un ELO mientras se está actualizando, eliminando las
  condiciones de carrera sobre la sección crítica del puntaje.

  \item \textbf{Una variable de condición (\texttt{g\_cv})} sobre la que los
  jugadores esperan cuando no encuentran rival o no hay tablero libre. Esto evita
  el \emph{busy-waiting}: el hilo se bloquea y solo es despertado cuando llega un
  candidato nuevo a la cola o cuando se libera un tablero.

  \item \textbf{Un contador de tableros libres} (en vez de un semáforo) protegido
  por \texttt{g\_mutex}. Se decrementa al iniciar una partida y se incrementa al
  terminarla, lo que impone el tope de $K$ partidas simultáneas usando solo
  primitivas de \texttt{pthread}.

  \item \textbf{Un mutex y una variable de condición propios de cada partida}
  (\texttt{m->mtx}, \texttt{m->cv}), que sincronizan al hilo del tablero con los
  dos jugadores por turnos. El tablero avisa a quién le toca, espera su jugada y
  cambia el turno; los jugadores esperan bloqueados su turno. Como tres hilos
  comparten la misma condición, usamos \texttt{pthread\_cond\_broadcast} y
  re-evaluamos la condición tras cada despertar para evitar despertares perdidos.
  Tener un candado por partida permite que las distintas partidas avancen en
  paralelo sin contención entre sí.

  \item \textbf{Una bandera \texttt{volatile sig\_atomic\_t}} para \texttt{SIGINT}:
  el manejador solo marca la bandera (única operación segura en ese contexto) y
  es el hilo principal quien, al detectarla, toma \texttt{g\_mutex}, deja de
  aceptar nuevas partidas y despierta con \texttt{broadcast} a todos los
  jugadores en espera para que terminen ordenadamente.
\end{itemize}

% ---------------------------------------------------------------------------
\section{Principales desafíos}

El mayor desafío fue \textbf{coordinar tres hilos en una misma partida} (dos
jugadores y el tablero) por turnos, sin caer en despertares perdidos ni en
interbloqueos. Lo resolvimos con un esquema en el que el tablero actúa como único
árbitro, todas las transiciones de estado se anuncian con \texttt{broadcast} y
cada hilo vuelve a comprobar la condición al despertar, de modo que el orden de
ejecución de los hilos no afecta la corrección.

Un segundo desafío fue el \textbf{manejo del ciclo de vida de la memoria de cada
partida}: liberar la estructura compartida sin que algún hilo la siga usando.
Para ello la partida lleva un contador de jugadores que ya salieron de ella; el
hilo del tablero espera a que ambos jugadores hayan terminado antes de liberar la
memoria, garantizando que nadie acceda a una región ya liberada.

También fue importante lograr un \textbf{cierre realmente ordenado}: tras recibir
\texttt{SIGINT} no deben quedar hilos colgados. Para esto el hilo principal espera
con \texttt{pthread\_join} a todos los jugadores y, además, a que el contador de
tableros activos llegue a cero antes de guardar el \emph{snapshot} y terminar.
Por último, respetar el tope de $K$ partidas \emph{sin busy-waiting} nos obligó a
combinar el contador de tableros con la variable de condición, en lugar de
recurrir a esperas activas.

% ---------------------------------------------------------------------------
\section{Conclusiones}

Durante el desarrollo de esta tarea, como grupo decidimos implementar
\texttt{cmatch} en C (C17), ya que es un lenguaje con el que nos sentimos cómodos
y que expone directamente las primitivas de \texttt{pthread}, lo que nos permitió
razonar con claridad sobre cada sección crítica. La principal decisión de diseño
fue separar la sincronización en dos niveles: un candado global para el estado
compartido de los jugadores y candados independientes por partida. Esta
separación mantiene la corrección del recurso más delicado --- el ELO --- y, al
mismo tiempo, permite que las partidas se jueguen en paralelo sin estorbarse.

Somos conscientes de que el uso de un único mutex global para la tabla de
jugadores es una solución de grano grueso que limita la concurrencia del
matchmaking; la preferimos por sobre un esquema de candados más finos porque
prioriza la corrección y es mucho más fácil de razonar, dejando la optimización
como un trabajo futuro. La implementación se validó con la herramienta
ThreadSanitizer, que no reportó condiciones de carrera, y se verificó que el
cierre por \texttt{SIGINT} y la restauración desde el \emph{snapshot} funcionan
correctamente. En conjunto, la tarea nos permitió aplicar de forma concreta los
conceptos de exclusión mutua, variables de condición y comunicación entre hilos
vistos en el curso.

% ---------------------------------------------------------------------------
\section*{Anexo}
\addcontentsline{toc}{section}{Anexo}
Repositorio con el código fuente, el \texttt{Makefile}, la configuración de
ejemplo y las instrucciones de uso:
\url{https://github.com/sephppx/Tarea_02-Sistema-Operativos}.

\end{document}
